\documentclass{article}
\usepackage[utf8]{inputenc}
\usepackage{mathtools}
\usepackage{amssymb}
\usepackage{centernot}
% New symbols
\let\oldsqrt\sqrt
\def\sqrt{\mathpalette\DHLhksqrt}
\def\DHLhksqrt#1#2{
\setbox0=\hbox{$#1\oldsqrt{#2\,}$}\dimen0=\ht0
\advance\dimen0-0.2\ht0
\setbox2=\hbox{\vrule height\ht0 depth -\dimen0}
{\box0\lower0.4pt\box2}}
\newcommand{\iu}{\mathrm{i}\mkern1mu}
\DeclarePairedDelimiter\abs{\lvert}{\rvert}
\DeclarePairedDelimiter\norm{\lVert}{\rVert}
\makeatletter
\let\oldabs\abs
\def\abs{\@ifstar{\oldabs}{\oldabs*}}
\let\oldnorm\norm
\def\norm{\@ifstar{\oldnorm}{\oldnorm*}}
\makeatother
\newcommand*{\Value}{\frac{1}{2}x^2}
\newcommand{\intab}{\int_a^b}
% End new symbols
\begin{document}

\section{\(\delta < 0, denominatore II grado\)}
\[\int \frac{1}{x^2 + 4x + 9} dx\]
Osserviamo che \(\int \frac{1}{x^2 + 1} dx = \arctan x + c\).
Provo allora a costruire qualcosa di simile all'arcotangente.
\[\int \frac{1}{x^2 + 4x + 9} dx = \int \frac{1}{x^2 + 4x + 4 + 5} dx = 5 \int \frac{1}{\frac{(x + 2)^2}{5} + 1} = 5 \arctan (\frac{x + 2}{\sqrt{5}}) + c\]

\end{document}